\documentclass[aspectratio=169,10pt]{beamer}

% Shared Beamer setup for Fall 2026 course slides.
\mode<presentation>{
  \usetheme{Madrid}
  \usecolortheme{default}
}

\definecolor{CalPolyGreen}{HTML}{154734}
\definecolor{CalPolyGold}{HTML}{BD8B13}
\definecolor{DeckInk}{HTML}{1F2933}
\definecolor{DeckMuted}{HTML}{5F6B7A}

\setbeamercolor{structure}{fg=CalPolyGreen}
\setbeamercolor{palette primary}{bg=CalPolyGreen,fg=white}
\setbeamercolor{palette secondary}{bg=DeckInk,fg=white}
\setbeamercolor{palette tertiary}{bg=CalPolyGold,fg=DeckInk}
\setbeamercolor{frametitle}{bg=CalPolyGreen,fg=white}
\setbeamercolor{title}{fg=CalPolyGreen}
\setbeamercolor{normal text}{fg=DeckInk,bg=white}
\setbeamercolor{block title}{bg=CalPolyGreen,fg=white}
\setbeamercolor{block body}{bg=black!3,fg=DeckInk}

\setbeamertemplate{navigation symbols}{}
\setbeamertemplate{footline}[frame number]
\setbeamertemplate{itemize item}{\color{CalPolyGold}$\blacktriangleright$}
\setbeamertemplate{itemize subitem}{\color{CalPolyGreen}$\bullet$}

\newcommand{\CourseNumber}{Course}
\newcommand{\CourseTitle}{Course Title}
\newcommand{\LectureNumber}{0}
\newcommand{\LectureTitle}{Lecture Title}
\newcommand{\LectureDate}{Date}
\newcommand{\CourseUnit}{Unit}

\newcommand{\coursenumber}[1]{\renewcommand{\CourseNumber}{#1}}
\newcommand{\coursetitle}[1]{\renewcommand{\CourseTitle}{#1}}
\newcommand{\lecturenumber}[1]{\renewcommand{\LectureNumber}{#1}}
\newcommand{\lecturetitle}[1]{\renewcommand{\LectureTitle}{#1}}
\newcommand{\lecturedate}[1]{\renewcommand{\LectureDate}{#1}}
\newcommand{\courseunit}[1]{\renewcommand{\CourseUnit}{#1}}

\author{Kirk Duran}
\institute{California Polytechnic State University, San Luis Obispo}
\date{\LectureDate}

\newcommand{\makedecktitle}{
  \begin{frame}[plain]
    \vspace{0.25in}
    {\usebeamerfont{title}\color{CalPolyGreen}\LectureTitle\par}
    \vspace{0.18in}
    {\large \CourseNumber{}: \CourseTitle\par}
    \vspace{0.08in}
    {\large Lecture \LectureNumber{} -- \LectureDate\par}
    \vfill
    {\small Kirk Duran \textbar{} Cal Poly, San Luis Obispo\par}
  \end{frame}
}

\newcommand{\placeholdernote}{
  \begin{block}{Placeholder}
    Replace these bullets with the final lecture material, examples, and in-class activities.
  \end{block}
}

\AtBeginSection[]{
  \begin{frame}{Roadmap}
    \tableofcontents[currentsection]
  \end{frame}
}


\pdfstringdefDisableCommands{\def\translate#1{#1}}

\graphicspath{{../assets/lecture-07/}}

% If an image file is not yet available, the slide displays a labeled
% placeholder using the intended filename. Place the matching file in
% ../assets/lecture-07/ to replace the placeholder automatically.
\newcommand{\lectureimage}[2][1.75in]{%
  \IfFileExists{../assets/lecture-07/#2}{%
    \includegraphics[width=\linewidth,height=#1,keepaspectratio]{#2}%
  }{%
    \fbox{%
      \parbox[c][#1][c]{0.94\linewidth}{%
        \centering
        \small\textbf{Image Placeholder}\\[0.08in]
        \scriptsize\texttt{\detokenize{#2}}
      }%
    }%
  }%
}

% Inline Python without repeatedly escaping underscores and operators.
\newcommand{\py}[1]{\texttt{\detokenize{#1}}}

\setbeamersize{text margin left=0.42in,text margin right=0.42in}

\coursenumber{CS 1001}
\coursetitle{Introduction to Computing}
\lecturenumber{7}
\lecturetitle{Code Tracing and Trace Tables}
\lecturedate{Fall 2026}
\courseunit{1. Computing and Program Execution}

\begin{document}

\makedecktitle

\begin{frame}{Today}
  \begin{itemize}
    \item How to reason about a Python program as a sequence of state changes.
    \item How to build a trace table using \textbf{Line \#} and \textbf{Known bindings}.
    \item How to distinguish program \textbf{state} from program \textbf{output}.
    \item How assignment, reassignment, arithmetic, strings, conversions, comparisons, and Boolean operators interact during execution.
    \item How to stop a trace at the first statement that cannot execute successfully.
  \end{itemize}

  \begin{keyidea}
    Tracing is not a new Python feature. It is a disciplined method for explaining what already-known Python statements do, one statement at a time.
  \end{keyidea}
\end{frame}

\begin{frame}{Scope: Straight-Line Programs Only}
  \begin{columns}[T,onlytextwidth]
    \begin{column}{0.47\textwidth}
      \begin{block}{Included}
        \begin{itemize}
          \item assignment and reassignment,
          \item arithmetic expressions,
          \item \py{int}, \py{float}, \py{str}, and \py{bool},
          \item type conversions,
          \item comparisons,
          \item \py{and}, \py{or}, and \py{not},
          \item \py{print()}.
        \end{itemize}
      \end{block}
    \end{column}

    \begin{column}{0.47\textwidth}
      \begin{alertblock}{Not yet}
        \begin{itemize}
          \item conditionals,
          \item loops,
          \item functions,
          \item lists,
          \item user input,
          \item exception handling,
          \item repeated or nested control flow.
        \end{itemize}
      \end{alertblock}
    \end{column}
  \end{columns}

  \begin{keyidea}
    Every program in this lecture executes exactly once from top to bottom.
  \end{keyidea}
\end{frame}

\sectionframe{Part 1: Programs as Sequences of State Changes}

\begin{frame}[shrink=8]{Execution Has an Order}
  \begin{columns}[T,onlytextwidth]
    \begin{column}{0.55\textwidth}
      \begin{itemize}
        \item Python executes a straight-line program in textual order: first statement, then second statement, then third statement, and so on.
        \item Each statement sees the program state produced by all successfully completed statements above it.
        \item A later assignment can change the state available to statements below it.
        \item A later assignment does \emph{not} travel backward and recompute earlier assignments.
      \end{itemize}

      \begin{block}{Execution order}
        \centering
        \texttt{statement 1} $\rightarrow$ \texttt{statement 2} $\rightarrow$ \texttt{statement 3} $\rightarrow$ $\cdots$
      \end{block}
    \end{column}

    \begin{column}{0.39\textwidth}
      % Search direction: vertical program lines with an execution arrow moving downward.
      \lectureimage[2.55in]{slide5.png}
    \end{column}
  \end{columns}
\end{frame}

\begin{frame}{Program State Means the Current Variable Associations}
  \begin{cpdefinition}
    The \textbf{program state} is the collection of variable names and the values currently associated with those names at a particular moment during execution.
  \end{cpdefinition}

  \begin{columns}[T,onlytextwidth]
    \begin{column}{0.47\textwidth}
      \begin{block}{After one statement}
        \py{x = 5}
        \vspace{0.10in}

        State: \py{x} $\rightarrow$ \py{5}
      \end{block}
    \end{column}

    \begin{column}{0.47\textwidth}
      \begin{block}{After another statement}
        \py{y = x + 2}
        \vspace{0.10in}

        State: \py{x} $\rightarrow$ \py{5}, \quad \py{y} $\rightarrow$ \py{7}
      \end{block}
    \end{column}
  \end{columns}

  \begin{keyidea}
    A trace table is a history of these states.
  \end{keyidea}
\end{frame}

\begin{frame}{Trace One Assignment}
  \begin{block}{Program}
    \centering
    \Large\py{x = 5}
  \end{block}

  \begin{center}
    \renewcommand{\arraystretch}{1.35}
    \begin{tabular}{c|l}
      \textbf{Line \#} & \textbf{Known bindings} \\\hline
      1 & \emph{none} \\
      End & \py{x -> 5}
    \end{tabular}
  \end{center}

  \begin{enumerate}
    \item Evaluate the right-hand side: \py{5} is already a value.
    \item Associate \py{x} with \py{5}.
    \item Advance to the next line; the new binding is visible in the next row of the trace.
  \end{enumerate}
\end{frame}

\begin{frame}{Dependent Assignments Use the Current State}
  \begin{block}{Program}
    \py{x = 5}\\[0.06in]
    \py{y = x * 2}
  \end{block}

  \begin{center}
    \renewcommand{\arraystretch}{1.30}
    \begin{tabular}{c|l}
      \textbf{Line \#} & \textbf{Known bindings} \\\hline
      1 & \emph{none} \\
      2 & \py{x -> 5} \\
      End & \py{x -> 5, y -> 10}
    \end{tabular}
  \end{center}

  \begin{keyidea}
    When the second statement executes, \py{x} already has a current value. Python retrieves that value before evaluating \py{x * 2}.
  \end{keyidea}
\end{frame}

\begin{frame}{Activity: Predict the State Before We Trace}
  \begin{block}{Program}
    \py{base = 8}\\
    \py{height = base + 4}\\
    \py{area = base * height}
  \end{block}

  \begin{activity}
    Without running Python, determine the bindings known when execution reaches each line. Include an \textbf{End} row to show the state after the final statement.
  \end{activity}

  \begin{block}{Questions to ask yourself}
    What values exist before each statement? Which names are read? Which name, if any, is updated?
  \end{block}
\end{frame}

\sectionframe{Part 2: A Disciplined Tracing Procedure}

\begin{frame}{The Five-Step Tracing Procedure}
  \begin{enumerate}
    \item Record the current line number and all bindings known \textbf{before} that line executes.
    \item Evaluate the statement using only those currently known bindings.
    \item If the statement is an assignment, update the set of bindings.
    \item If the statement is \py{print()}, record the displayed output without changing the bindings.
    \item Advance to the next line and record the updated bindings there. After the final line, record an \textbf{End} row.
  \end{enumerate}

  \begin{keyidea}
    A numbered row describes the state \emph{when execution reaches that line}. The effects of that line appear in the next row.
  \end{keyidea}
\end{frame}

\begin{frame}[shrink=8]{A Trace Table Separates Time, State, and Output}
  \begin{columns}[T,onlytextwidth]
    \begin{column}{0.54\textwidth}
      \begin{itemize}
        \item The \textbf{Line \#} column identifies the statement Python is about to execute.
        \item \textbf{Known bindings} records every variable-value association already established when execution reaches that line.
        \item Therefore, line 1 begins with no known program-variable bindings.
        \item The effect of an assignment first appears on the \emph{next} row; an \textbf{End} row records the state after the final statement.
        \item Program output is recorded separately so that bindings and displayed output are not conflated.
      \end{itemize}
    \end{column}

    \begin{column}{0.40\textwidth}
      % Search direction: annotated trace table showing statement, state, and output regions.
      \lectureimage[2.45in]{slide12.png}
    \end{column}
  \end{columns}
\end{frame}

\begin{frame}{Worked Trace: Assignment, Dependency, Reassignment, Output}
  \begin{block}{Program}
    \py{x = 5}\\
    \py{y = x * 2}\\
    \py{x = y - 3}\\
    \py{print(x)}
  \end{block}

  \begin{center}
    \scriptsize
    \renewcommand{\arraystretch}{1.25}
    \begin{tabular}{c|l}
      \textbf{Line \#} & \textbf{Known bindings} \\\hline
      1 & \emph{none} \\
      2 & \py{x -> 5} \\
      3 & \py{x -> 5, y -> 10} \\
      4 & \py{x -> 7, y -> 10} \\
      End & \py{x -> 7, y -> 10}
    \end{tabular}

    \vspace{0.10in}
    \textbf{Output produced at line 4:} \py{7}
  \end{center}
\end{frame}

\begin{frame}{Reassignment Replaces a Variable's Current Association}
  \begin{columns}[T,onlytextwidth]
    \begin{column}{0.52\textwidth}
      \begin{block}{Program fragment}
        \py{x = 5}\\
        \py{x = 10}
      \end{block}

      \begin{itemize}
        \item After the first statement, \py{x} refers to \py{5}.
        \item After the second statement, \py{x} refers to \py{10}.
        \item The old association is no longer the current state for \py{x}.
      \end{itemize}
    \end{column}

    \begin{column}{0.42\textwidth}
      % Search direction: variable x first pointing to 5, then arrow redirected to 10.
      \lectureimage[2.35in]{slide14.png}
    \end{column}
  \end{columns}
\end{frame}

\begin{frame}[shrink=8]{Reassignment Does Not Recompute Earlier Assignments}
  \begin{block}{Program}
    \py{x = 3}\\
    \py{y = x + 2}\\
    \py{x = 10}
  \end{block}

  \begin{center}
    \renewcommand{\arraystretch}{1.28}
    \begin{tabular}{c|l}
      \textbf{Line \#} & \textbf{Known bindings} \\\hline
      1 & \emph{none} \\
      2 & \py{x -> 3} \\
      3 & \py{x -> 3, y -> 5} \\
      End & \py{x -> 10, y -> 5}
    \end{tabular}
  \end{center}

  \begin{keyidea}
    Assignment stores the value computed at that moment. \py{y} does not contain a live formula that automatically follows future changes to \py{x}.
  \end{keyidea}
\end{frame}

\begin{frame}{Activity: Trace Reassignment Carefully}
  \begin{block}{Program}
    \py{a = 4}\\
    \py{b = a + 6}\\
    \py{a = b * 2}\\
    \py{c = a - b}
  \end{block}

  \begin{activity}
    Build the trace table using \textbf{Line \#} and \textbf{Known bindings}. Circle the row where the binding for \py{a} changes from its earlier value.
  \end{activity}
\end{frame}

\sectionframe{Part 3: State Is Not Output}

\begin{frame}{Assignment Changes State; \texttt{print()} Produces Output}
  \begin{columns}[T,onlytextwidth]
    \begin{column}{0.47\textwidth}
      \begin{block}{Assignment}
        \centering
        \Large\py{x = 7}
      \end{block}
      \begin{itemize}
        \item Updates program state.
        \item Displays nothing by itself.
      \end{itemize}
    \end{column}

    \begin{column}{0.47\textwidth}
      \begin{block}{Output}
        \centering
        \Large\py{print(x)}
      \end{block}
      \begin{itemize}
        \item Reads the current value of \py{x}.
        \item Displays that value.
        \item Does not reassign \py{x}.
      \end{itemize}
    \end{column}
  \end{columns}
\end{frame}

\begin{frame}{Trace State and Record Output Separately}
  \begin{block}{Program}
    \py{price = 12}\\
    \py{quantity = 3}\\
    \py{total = price * quantity}\\
    \py{print(total)}
  \end{block}

  \begin{center}
    \scriptsize
    \renewcommand{\arraystretch}{1.22}
    \begin{tabular}{c|l}
      \textbf{Line \#} & \textbf{Known bindings} \\\hline
      1 & \emph{none} \\
      2 & \py{price -> 12} \\
      3 & \py{price -> 12, quantity -> 3} \\
      4 & \py{price -> 12, quantity -> 3, total -> 36} \\
      End & \py{price -> 12, quantity -> 3, total -> 36}
    \end{tabular}

    \vspace{0.10in}
    \textbf{Output produced at line 4:} \py{36}
  \end{center}
\end{frame}

\begin{frame}{Activity: Predict Both State and Output}
  \begin{block}{Program}
    \py{minutes = 90}\\
    \py{hours = minutes / 60}\\
    \py{print(hours)}\\
    \py{minutes = 120}
  \end{block}

  \begin{activity}
    Trace the program. After the final statement, what is the state? What was displayed earlier? Keep those answers separate.
  \end{activity}
\end{frame}

\sectionframe{Part 4: Tracing Arithmetic and Basic Types}

\begin{frame}{Arithmetic Expressions Are Evaluated Before Assignment}
  \begin{block}{Program}
    \py{width = 8}\\
    \py{height = 5}\\
    \py{area = width * height}\\
    \py{perimeter = 2 * (width + height)}
  \end{block}

  \begin{block}{Evaluation of one statement}
    \centering
    \py{perimeter = 2 * (width + height)}\\[0.08in]
    $\downarrow$ substitute current values\\[0.08in]
    \py{perimeter = 2 * (8 + 5)} $\rightarrow$ \py{perimeter = 26}
  \end{block}
\end{frame}

\begin{frame}{Code Writing: Rectangle Calculator}
  \begin{activity}
    Write a straight-line Python script for a rectangle with width \py{7.5} and height \py{4}.

    Your script should:
    \begin{enumerate}
      \item store the width and height,
      \item calculate the area,
      \item calculate the perimeter,
      \item print the area,
      \item print the perimeter.
    \end{enumerate}
  \end{activity}

  \begin{keyidea}
    Before running it, predict the value and type of every variable.
  \end{keyidea}
\end{frame}

\begin{frame}{One Possible Rectangle Calculator}
  \begin{block}{Script}
    \py{width = 7.5}\\
    \py{height = 4}\\
    \py{area = width * height}\\
    \py{perimeter = 2 * (width + height)}\\
    \py{print(area)}\\
    \py{print(perimeter)}
  \end{block}

  \begin{itemize}
    \item \py{width} is a \py{float}; \py{height} is an \py{int}.
    \item The arithmetic expressions involving \py{width} produce floating-point results.
    \item The two \py{print()} statements produce output but do not modify the four variables.
  \end{itemize}
\end{frame}

\begin{frame}[shrink=8]{Division Often Introduces a Floating-Point Value}
  \begin{block}{Program}
    \py{distance = 150}\\
    \py{time = 2}\\
    \py{speed = distance / time}
  \end{block}

  \begin{center}
    \renewcommand{\arraystretch}{1.30}
    \begin{tabular}{c|l}
      \textbf{Line \#} & \textbf{Known bindings} \\\hline
      1 & \emph{none} \\
      2 & \py{distance -> 150} \\
      3 & \py{distance -> 150, time -> 2} \\
      End & \py{distance -> 150, time -> 2, speed -> 75.0}
    \end{tabular}
  \end{center}

  \begin{keyidea}
    A correct trace records both value and, when relevant, type: \py{75.0} is not written as \py{75}.
  \end{keyidea}
\end{frame}

\begin{frame}{Strings Can Depend on Numeric Computations}
  \begin{block}{Program}
    \py{units = 4}\\
    \py{hours = 3.5}\\
    \py{workload = units * hours}\\
    \py{label = "Estimated workload: " + str(workload)}
  \end{block}

  \begin{block}{Last statement}
    \py{workload} has the numeric value \py{14.0}.\\[0.08in]
    \py{str(workload)} produces the string \py{"14.0"}.\\[0.08in]
    Concatenation produces \py{"Estimated workload: 14.0"}.
  \end{block}
\end{frame}

\begin{frame}{Trace Mixed Numeric and String Values}
  \begin{center}
    \scriptsize
    \renewcommand{\arraystretch}{1.20}
    \begin{tabular}{c|l}
      \textbf{Line \#} & \textbf{Known bindings} \\\hline
      1 & \emph{none} \\
      2 & \py{units -> 4} \\
      3 & \py{units -> 4, hours -> 3.5} \\
      4 & \py{units -> 4, hours -> 3.5, workload -> 14.0} \\
      End & \py{units -> 4, hours -> 3.5, workload -> 14.0, label -> "Estimated workload: 14.0"}
    \end{tabular}
  \end{center}

  \begin{keyidea}
    The trace records the value actually stored in each variable, not merely the source expression that produced it.
  \end{keyidea}
\end{frame}

\begin{frame}{Code Writing: Workload Estimate}
  \begin{activity}
    Write a straight-line script that stores:
    \begin{itemize}
      \item \py{courses = 5},
      \item \py{hours_per_course = 4.5},
      \item the total weekly hours,
      \item a string such as \py{"Weekly hours: 22.5"}.
    \end{itemize}

    Print the final string. Use \py{str()} only where it is needed.
  \end{activity}
\end{frame}

\begin{frame}{One Possible Workload Script}
  \begin{block}{Script}
    \py{courses = 5}\\
    \py{hours_per_course = 4.5}\\
    \py{weekly_hours = courses * hours_per_course}\\
    \py{summary = "Weekly hours: " + str(weekly_hours)}\\
    \py{print(summary)}
  \end{block}

  \begin{itemize}
    \item The arithmetic result is \py{22.5}.
    \item The conversion produces \py{"22.5"} for concatenation.
    \item The variable \py{weekly_hours} remains numeric; converting it for the string does not reassign it.
  \end{itemize}
\end{frame}

\sectionframe{Part 5: Tracing Comparisons and Boolean Values}

\begin{frame}{Comparisons Produce Values That Can Be Traced}
  \begin{block}{Program}
    \py{score = 84}\\
    \py{minimum_score = 70}\\
    \py{is_passing = score >= minimum_score}
  \end{block}

  \begin{block}{Evaluation}
    \centering
    \py{score >= minimum_score} $\rightarrow$ \py{84 >= 70} $\rightarrow$ \py{True}
  \end{block}

  \begin{keyidea}
    The comparison is evaluated first; assignment then associates \py{is_passing} with the Boolean result.
  \end{keyidea}
\end{frame}

\begin{frame}{Boolean Variables Become Part of Program State}
  \begin{block}{Program}
    \py{score = 84}\\
    \py{minimum_score = 70}\\
    \py{is_passing = score >= minimum_score}\\
    \py{is_honors = score >= 90}\\
    \py{eligible = is_passing and not is_honors}
  \end{block}

  \begin{center}
    \scriptsize
    \renewcommand{\arraystretch}{1.16}
    \begin{tabular}{c|l}
      \textbf{Line \#} & \textbf{Known bindings} \\\hline
      1 & \emph{none} \\
      2 & \py{score -> 84} \\
      3 & \py{score -> 84, minimum_score -> 70} \\
      4 & \py{score -> 84, minimum_score -> 70, is_passing -> True} \\
      5 & \py{score -> 84, minimum_score -> 70, is_passing -> True, is_honors -> False} \\
      End & \py{score -> 84, minimum_score -> 70, is_passing -> True, is_honors -> False, eligible -> True}
    \end{tabular}
  \end{center}
\end{frame}

\begin{frame}{Do Not Confuse Assignment with Equality Comparison}
  \begin{columns}[T,onlytextwidth]
    \begin{column}{0.47\textwidth}
      \begin{block}{Assignment}
        \centering
        \Large\py{x = 10}
      \end{block}
      Changes the state of \py{x}.
    \end{column}

    \begin{column}{0.47\textwidth}
      \begin{block}{Comparison}
        \centering
        \Large\py{x == 10}
      \end{block}
      Computes \py{True} or \py{False}; does not change \py{x}.
    \end{column}
  \end{columns}

  \begin{activity}
    Suppose \py{x} currently equals \py{7}. Describe the effect and result of each expression above.
  \end{activity}
\end{frame}

\begin{frame}{Code Writing: A Simple Budget Check}
  \begin{activity}
    Write a straight-line script that stores:
    \begin{itemize}
      \item \py{budget = 100},
      \item \py{item_price = 64.50},
      \item \py{tax_rate = 0.08},
      \item the tax,
      \item the total cost,
      \item a Boolean named \py{within_budget} that compares the total cost with the budget.
    \end{itemize}

    Print \py{total_cost} and \py{within_budget}. Do not use a conditional statement.
  \end{activity}
\end{frame}

\begin{frame}{One Possible Budget Script}
  \begin{block}{Script}
    \py{budget = 100}\\
    \py{item_price = 64.50}\\
    \py{tax_rate = 0.08}\\
    \py{tax = item_price * tax_rate}\\
    \py{total_cost = item_price + tax}\\
    \py{within_budget = total_cost <= budget}\\
    \py{print(total_cost)}\\
    \py{print(within_budget)}
  \end{block}

  \begin{keyidea}
    A Boolean can summarize a relationship among previously computed values even though the program does not yet make a decision based on it.
  \end{keyidea}
\end{frame}

\sectionframe{Part 6: Full Trace Synthesis}

\begin{frame}{Integrated Example}
  \begin{block}{Program}
    \py{x = 3}\\
    \py{y = x + 2}\\
    \py{x = 10}\\
    \py{z = x + y}\\
    \py{larger = z > 10}\\
    \py{same = x == y}\\
    \py{result = larger and not same}\\
    \py{print(result)}
  \end{block}

  \begin{activity}
    Before tracing, predict the final value and type of every variable.
  \end{activity}
\end{frame}

\begin{frame}[shrink=8]{Integrated Example: Complete Trace}
  \begin{center}
    \scriptsize
    \renewcommand{\arraystretch}{1.16}
    \begin{tabular}{c|l}
      \textbf{Line \#} & \textbf{Known bindings} \\\hline
      1 & \emph{none} \\
      2 & \py{x -> 3} \\
      3 & \py{x -> 3, y -> 5} \\
      4 & \py{x -> 10, y -> 5} \\
      5 & \py{x -> 10, y -> 5, z -> 15} \\
      6 & \py{x -> 10, y -> 5, z -> 15, larger -> True} \\
      7 & \py{x -> 10, y -> 5, z -> 15, larger -> True, same -> False} \\
      8 & \py{x -> 10, y -> 5, z -> 15, larger -> True, same -> False, result -> True} \\
      End & \py{x -> 10, y -> 5, z -> 15, larger -> True, same -> False, result -> True}
    \end{tabular}

    \vspace{0.08in}
    \textbf{Output produced at line 8:} \py{True}
  \end{center}
\end{frame}

\begin{frame}{Read the Integrated Trace as a Sequence of Questions}
  \begin{enumerate}
    \item When \py{y = x + 2} executes, which \py{x} exists? \textbf{3}.
    \item After \py{x = 10}, does \py{y} change? \textbf{No}; it remains \textbf{5}.
    \item What values does \py{z = x + y} read? \textbf{10} and \textbf{5}.
    \item What does \py{z > 10} compare? \textbf{15} and \textbf{10}.
    \item What does \py{x == y} compare? \textbf{10} and \textbf{5}.
    \item What does \py{larger and not same} evaluate to? \textbf{True}.
    \item What changes on the final line? \textbf{Output only}.
  \end{enumerate}
\end{frame}

\begin{frame}{Activity: Trace a Calculation Script}
  \begin{block}{Program}
    \py{miles = 180}\\
    \py{gallons = 6}\\
    \py{mpg = miles / gallons}\\
    \py{target = 28}\\
    \py{meets_target = mpg >= target}\\
    \py{message = "MPG: " + str(mpg)}\\
    \py{print(message)}\\
    \py{print(meets_target)}
  \end{block}

  \begin{activity}
    Construct the complete trace table using Line \# and Known bindings. Include variable types in annotations beside your table.
  \end{activity}
\end{frame}

\sectionframe{Part 7: Tracing to the First Error}

\begin{frame}{A Trace Stops When a Statement Cannot Complete}
  \begin{block}{Program}
    \py{count = 5}\\
    \py{label = "Count: " + count}\\
    \py{double = count * 2}
  \end{block}

  \begin{itemize}
    \item The first statement executes successfully: \py{count} becomes \py{5}.
    \item The second statement attempts to concatenate a string and an integer directly.
    \item That statement does not successfully produce a value for \py{label}.
    \item Execution stops there; the later assignment to \py{double} is not traced as executed.
  \end{itemize}

  \begin{alertblock}{Tracing rule}
    Identify the first failing statement. Do not invent state changes from statements that were never reached.
  \end{alertblock}
\end{frame}

\begin{frame}{Trace Only the Successfully Completed State}
  \begin{center}
    \renewcommand{\arraystretch}{1.28}
    \begin{tabular}{c|l}
      \textbf{Line \#} & \textbf{Known bindings} \\\hline
      1 & \emph{none} \\
      2 & \py{count -> 5}
    \end{tabular}

    \vspace{0.10in}
    \textbf{Line 2 fails.} No binding for \py{label} is created; execution stops with \py{count -> 5}.
  \end{center}

  \begin{activity}
    Rewrite only the failing statement so that the program can continue, using a type conversion already introduced in this course.
  \end{activity}
\end{frame}

\begin{frame}{Common Tracing Errors}
  \begin{itemize}
    \item Reading the whole program as though all assignments happen simultaneously.
    \item Forgetting that reassignment replaces a variable's current value.
    \item Recomputing a previous assignment after one of its source variables changes.
    \item Treating \py{print()} as though it changes program state.
    \item Confusing \py{=} with \py{==}.
    \item Treating \py{"5"} and \py{5} as interchangeable values.
    \item Skipping rows and losing the exact moment when state changes.
    \item Reporting only the final answer instead of showing the execution history.
  \end{itemize}
\end{frame}

\begin{frame}{A Useful Self-Check for Every Row}
  \begin{columns}[T,onlytextwidth]
    \begin{column}{0.52\textwidth}
      \begin{enumerate}
        \item What statement just executed?
        \item What values did it read?
        \item What expression result did it compute?
        \item What variable, if any, changed?
        \item What output, if any, appeared?
      \end{enumerate}
    \end{column}

    \begin{column}{0.42\textwidth}
      % Search direction: checklist beside a small trace table, emphasizing one row at a time.
      \lectureimage[2.35in]{trace-row-checklist.png}
    \end{column}
  \end{columns}
\end{frame}

\sectionframe{Part 8: Writing Small Scripts You Can Trace}

\begin{frame}{Code Writing: Temperature Conversion}
  \begin{activity}
    Write a straight-line script that:
    \begin{enumerate}
      \item stores \py{fahrenheit = 77},
      \item computes Celsius using \py{(fahrenheit - 32) * 5 / 9},
      \item builds a string beginning with \py{"Celsius: "},
      \item prints the string,
      \item stores a Boolean named \py{is_warm} that is true when Celsius is at least \py{20}.
    \end{enumerate}

    Predict the known bindings at each line before running it, and include the final bindings at \textbf{End}.
  \end{activity}
\end{frame}

\begin{frame}{One Possible Temperature Script}
  \begin{block}{Script}
    \py{fahrenheit = 77}\\
    \py{celsius = (fahrenheit - 32) * 5 / 9}\\
    \py{label = "Celsius: " + str(celsius)}\\
    \py{print(label)}\\
    \py{is_warm = celsius >= 20}
  \end{block}

  \begin{itemize}
    \item The calculation produces \py{25.0}.
    \item \py{label} stores \py{"Celsius: 25.0"}.
    \item Printing the label does not change it.
    \item \py{is_warm} becomes \py{True}.
  \end{itemize}
\end{frame}

\begin{frame}{Code Writing: Purchase Summary}
  \begin{activity}
    Write a script for a purchase of \py{3} notebooks at \py{2.75} dollars each.

    Include variables for:
    \begin{itemize}
      \item quantity,
      \item unit price,
      \item subtotal,
      \item a fixed shipping cost of \py{4.00},
      \item final total,
      \item a Boolean that records whether the final total is less than \py{15},
      \item a printable string summary of the final total.
    \end{itemize}
  \end{activity}
\end{frame}

\begin{frame}{From Writing to Tracing}
  \begin{columns}[T,onlytextwidth]
    \begin{column}{0.48\textwidth}
      \begin{block}{When writing}
        \begin{itemize}
          \item Choose meaningful variable names.
          \item Build calculations from previously available values.
          \item Keep each statement conceptually simple.
          \item Predict values before running.
        \end{itemize}
      \end{block}
    \end{column}

    \begin{column}{0.48\textwidth}
      \begin{block}{When tracing}
        \begin{itemize}
          \item Respect top-to-bottom order.
          \item Use only the current state.
          \item Record every state-changing statement.
          \item Separate state from displayed output.
        \end{itemize}
      \end{block}
    \end{column}
  \end{columns}

  \begin{keyidea}
    Writing and tracing reinforce the same execution model from opposite directions.
  \end{keyidea}
\end{frame}

\end{document}
